% GENERATED FILE. Edit canonical lessons, then run npm run generate:books.
% canonical-source-hash: fnv1a64:5b479e5c9f52cd38
% canonical-lessons: RU-C296-spasatel, RU-C296-spasat, RU-C296-spaseniye, RU-C296-evakuirovat, RU-C296-evakuatsiya

\chapter{Rescue}
\label{ch:ru-rescue}
\hlchaptermodality{296}

\begin{chapteropening}
\textbf{By the end of this chapter:} \emph{I can talk about a rescuer, saving someone, a rescue, evacuating people and an evacuation.}

Complete the last lesson of chapter 296: I can talk about a rescuer, saving someone, a rescue, evacuating people and an evacuation.
\end{chapteropening}

\section[\texorpdfstring{\ru{спасатель} (spasátel)}{спасатель (spasátel)}]{\ru{спасатель} (spasátel) — a rescuer, a lifeguard}

\label{lesson:RU-C296-spasatel}



\begin{quote}
Before the new one: say the Russian for to catch fire, then the Russian for a fire extinguisher.
\end{quote}

\subsection*{You'll want to know: \ru{спасатель}}
\textbf{\ru{спасатель}} (\emph{spasátel}) — \textquotedblleft{}a rescuer, a lifeguard\textquotedblright{}.

\begin{quote}
\ru{Спасатели} \ru{приехали} \ru{через} \ru{десять} \ru{минут}.
\end{quote}

\begin{quote}
\emph{Spasáteli priyékhali chérez désyat minút.}
\end{quote}

\begin{quote}
The rescuers arrived ten minutes later.
\end{quote}

Genitive \emph{\ru{спасателя}}.

\begin{cousinweb}
\emph{\ru{спасатель}} is built on \emph{\ru{спасать}} \textquotedblleft{}to save\textquotedblright{}, from \emph{\ru{спасти}}, Old East Slavic \emph{sŭpasti} (EtymDB), the verb also behind \emph{\ru{спасибо}}, with the ending \emph{-\ru{тель}} for a person.
\end{cousinweb}

\begin{grammarlens}[title={What you've built}]
One more word for telling, explaining and describing at B1.
\end{grammarlens}

\subsection*{Your turn}
\begin{itemize}
\raggedright
  \item \emph{Say it:} \emph{spasátel}
  \item \emph{Say it:} \emph{spasátel}, once more
  \item \emph{Say it:} a sentence of your own with \emph{spasátel}
  \item \emph{Recall:} say the Russian for to catch fire, then the Russian for a fire extinguisher, then say \emph{spasátel} again
\end{itemize}

\begin{tcolorbox}[breakable,colback=teal!4,colframe=teal!35!black,title={Before you move on}]
What does \ru{спасатель} mean? (\textquotedblleft{}A rescuer, a lifeguard\textquotedblright{}.) Say it once more.
\end{tcolorbox}
\section[\texorpdfstring{\ru{спасать} (spasát)}{спасать (spasát)}]{\ru{спасать} (spasát) — to save, to rescue}

\label{lesson:RU-C296-spasat}



\begin{quote}
Before the new one: say the Russian for a fire extinguisher, then the Russian for a rescuer.
\end{quote}

\subsection*{You'll want to know: \ru{спасать}}
\textbf{\ru{спасать}} (\emph{spasát}) — \textquotedblleft{}to save, to rescue\textquotedblright{}.

\begin{quote}
\ru{Пожарный} \ru{спас} \ru{ребёнка}.
\end{quote}

\begin{quote}
\emph{Pozhárnyy spas rebyónka.}
\end{quote}

\begin{quote}
The firefighter saved a child.
\end{quote}

Imperfective; perfective \emph{\ru{спасти}} (\emph{spastí}), past \emph{\ru{спас}}, \emph{\ru{спасла}}. \emph{\ru{Спасите}!} — help!

\begin{cousinweb}
\emph{\ru{спасать}} is the imperfective of \emph{\ru{спасти}}, inherited through Old East Slavic \emph{sŭpasti}, from \emph{sŭ-} \textquotedblleft{}with\textquotedblright{} and \emph{pasti} \textquotedblleft{}to pasture, to herd\textquotedblright{}, PIE \emph{peh\textsubscript{2}-} \textquotedblleft{}to guard, to feed\textquotedblright{} (EtymDB); the same verb stands behind \emph{\ru{спасибо}}.
\end{cousinweb}

\begin{grammarlens}[title={What you've built}]
One more word for telling, explaining and describing at B1.
\end{grammarlens}

\subsection*{Your turn}
\begin{itemize}
\raggedright
  \item \emph{Say it:} \emph{spasát}
  \item \emph{Say it:} \emph{spasát}, once more
  \item \emph{Say it:} a sentence of your own with \emph{spasát}
  \item \emph{Recall:} say the Russian for a fire extinguisher, then the Russian for a rescuer, then say \emph{spasát} again
\end{itemize}

\begin{tcolorbox}[breakable,colback=teal!4,colframe=teal!35!black,title={Before you move on}]
What does \ru{спасать} mean? (\textquotedblleft{}To save, to rescue\textquotedblright{}.) Say it once more.
\end{tcolorbox}
\section[\texorpdfstring{\ru{спасение} (spaséniye)}{спасение (spaséniye)}]{\ru{спасение} (spaséniye) — rescue, salvation}

\label{lesson:RU-C296-spaseniye}



\begin{quote}
Before the new one: say the Russian for a rescuer, then the Russian for to save.
\end{quote}

\subsection*{You'll want to know: \ru{спасение}}
\textbf{\ru{спасение}} (\emph{spaséniye}) — \textquotedblleft{}rescue, salvation\textquotedblright{}.

\begin{quote}
\ru{Спасение} \ru{пришло} \ru{в} \ru{последнюю} \ru{минуту}.
\end{quote}

\begin{quote}
\emph{Spaséniye prishló v poslédnyuyu minútu.}
\end{quote}

\begin{quote}
Rescue came at the last minute.
\end{quote}

Neuter.

\begin{cousinweb}
\emph{\ru{спасение}} is built on \emph{\ru{спасти}} \textquotedblleft{}to save\textquotedblright{}, the perfective of \emph{\ru{спасать}}, Old East Slavic \emph{sŭpasti}, from \emph{pasti} \textquotedblleft{}to herd, to guard\textquotedblright{} (EtymDB).
\end{cousinweb}

\begin{grammarlens}[title={What you've built}]
One more word for telling, explaining and describing at B1.
\end{grammarlens}

\subsection*{Your turn}
\begin{itemize}
\raggedright
  \item \emph{Say it:} \emph{spaséniye}
  \item \emph{Say it:} \emph{spaséniye}, once more
  \item \emph{Say it:} a sentence of your own with \emph{spaséniye}
  \item \emph{Recall:} say the Russian for a rescuer, then the Russian for to save, then say \emph{spaséniye} again
\end{itemize}

\begin{tcolorbox}[breakable,colback=teal!4,colframe=teal!35!black,title={Before you move on}]
What does \ru{спасение} mean? (\textquotedblleft{}Rescue, salvation\textquotedblright{}.) Say it once more.
\end{tcolorbox}
\section[\texorpdfstring{\ru{эвакуировать} (evakuírovat)}{эвакуировать (evakuírovat)}]{\ru{эвакуировать} (evakuírovat) — to evacuate}

\label{lesson:RU-C296-evakuirovat}



\begin{quote}
Before the new one: say the Russian for to save, then the Russian for rescue.
\end{quote}

\subsection*{You'll want to know: \ru{эвакуировать}}
\textbf{\ru{эвакуировать}} (\emph{evakuírovat}) — \textquotedblleft{}to evacuate\textquotedblright{}.

\begin{quote}
\ru{Жителей} \ru{эвакуировали} \ru{из} \ru{дома}.
\end{quote}

\begin{quote}
\emph{Zhíteley evakuírovali iz dóma.}
\end{quote}

\begin{quote}
The residents were evacuated from the building.
\end{quote}

One form for both aspects. It takes the accusative.

\begin{cousinweb}
\emph{\ru{эвакуировать}} goes back to Latin \emph{ēvacuāre} \textquotedblleft{}to empty out\textquotedblright{}, built on \emph{vacuus} \textquotedblleft{}empty\textquotedblright{} (Lewis \& Short); EtymDB lists the Russian verb with no parent.
\end{cousinweb}

\begin{grammarlens}[title={What you've built}]
One more word for telling, explaining and describing at B1.
\end{grammarlens}

\subsection*{Your turn}
\begin{itemize}
\raggedright
  \item \emph{Say it:} \emph{evakuírovat}
  \item \emph{Say it:} \emph{evakuírovat}, once more
  \item \emph{Say it:} a sentence of your own with \emph{evakuírovat}
  \item \emph{Recall:} say the Russian for to save, then the Russian for rescue, then say \emph{evakuírovat} again
\end{itemize}

\begin{tcolorbox}[breakable,colback=teal!4,colframe=teal!35!black,title={Before you move on}]
What does \ru{эвакуировать} mean? (\textquotedblleft{}To evacuate\textquotedblright{}.) Say it once more.
\end{tcolorbox}
\section[\texorpdfstring{\ru{эвакуация} (evakuátsiya)}{эвакуация (evakuátsiya)}]{\ru{эвакуация} (evakuátsiya) — an evacuation}

\label{lesson:RU-C296-evakuatsiya}



\begin{quote}
Before the new one: say the Russian for rescue, then the Russian for to evacuate.
\end{quote}

\subsection*{You'll want to know: \ru{эвакуация}}
\textbf{\ru{эвакуация}} (\emph{evakuátsiya}) — \textquotedblleft{}an evacuation\textquotedblright{}.

\begin{quote}
\ru{Эвакуация} \ru{началась} \ru{в} \ru{пять} \ru{утра}.
\end{quote}

\begin{quote}
\emph{Evakuátsiya nachalás v pyat utrá.}
\end{quote}

\begin{quote}
The evacuation began at five in the morning.
\end{quote}

Feminine.

\begin{cousinweb}
\emph{\ru{эвакуация}} goes back to Latin \emph{ēvacuātiō} \textquotedblleft{}an emptying\textquotedblright{}, from \emph{ēvacuāre} and \emph{vacuus} \textquotedblleft{}empty\textquotedblright{} (Lewis \& Short); EtymDB lists no parent.
\end{cousinweb}

\begin{grammarlens}[title={What you've built}]
One more word for telling, explaining and describing at B1.
\end{grammarlens}

\subsection*{Your turn}
\begin{itemize}
\raggedright
  \item \emph{Say it:} \emph{evakuátsiya}
  \item \emph{Say it:} \emph{evakuátsiya}, once more
  \item \emph{Say it:} a sentence of your own with \emph{evakuátsiya}
  \item \emph{Recall:} say the Russian for rescue, then the Russian for to evacuate, then say \emph{evakuátsiya} again
\end{itemize}

\begin{tcolorbox}[breakable,colback=teal!4,colframe=teal!35!black,title={Before you move on}]
What does \ru{эвакуация} mean? (\textquotedblleft{}An evacuation\textquotedblright{}.) Say it once more.
\end{tcolorbox}
